% Nouveau EGA V (cote25): archive cote25, site122 (PDF page123).
% English translation of the checked French annotated master.
% AI-translated by GPT-6 Astra (OpenAI, ChatGPT Pro mode).
\scanpage{25/122}
\mspage{25}

linear systems of divisors without fixed component on X parametrized
by $\overset{\hand{\vee}}{P}$ , and the set $\hand{\unsure{\psi|\emptyset}}$ of rational maps from X into P, \hand{which}
\var{\inserted{over $\bar{k}$ factor through no hyperplane of $\unsure{\overline{P}|P}_{\bar{k}}$.}}{over $\bar{k}$ factor through no hyperplane of $\unsure{\overline{P}|P}_{\bar{k}}$.}{insertion}{ega25_25_122-v0001}

Indeed, \var{\inserted[typed]{when the local rings of X are factorial.}}{when the local rings of X are factorial.}{insertion}{ega25_25_122-v0002} \var{\deleted{\var{\xout{Since X is regular, \unsure{it|X} is the same for \unsure{$X\times_S P$}, which is smooth}}{}{typewriter strike-out}{ega25_25_122-v0004}
\var{\xout{over X, hence by a result of Auslander-Buchsbaum (which will be taken up again}}{}{typewriter strike-out}{ega25_25_122-v0005}
\var{\xout{in Chap III) \illegible{1} of the local rings of \unsure{$X\times_S P$} are factorial,}}{}{typewriter strike-out}{ega25_25_122-v0006}
and since codim(X-U,X) \quad 2}}{}{deletion}{ega25_25_122-v0003} , it follows that every invertible Module on
U can be extended to an invertible Module on X, so the condition
considered in b) is automatically satisfied. On the other hand, by
\var{\deleted{\unsure{a Memoir} \illegible{1} \unsure{in a} \unsure{later} no.}}{}{deletion}{ega25_25_122-v0007}
Auslander-Buchsbaum (\var{\deleted{which will appear one day in Chap III}}{}{deletion}{ega25_25_122-v0008}\var{\inserted{cf §20}}{cf §20}{insertion}{ega25_25_122-v0009}) a regular local
rin\hand{/}g is factorial.

\var{\margin{\unsure{$\times\quad\times\quad\nearrow$}}}{\placedmargin{\unsure{$\times\quad\times\quad\nearrow$}}}{marginal note placed}{ega25_25_122-v0010}Proposition 2 b) shows us that with an additional hypothesis
\var{\inserted[typed]{(namely $D\quad\hand{\unsure{\phi|\emptyset}}$)}}{(namely $D\quad\hand{\unsure{\phi|\emptyset}}$)}{insertion}{ega25_25_122-v0011}
on a\var{\xout{family}}{}{typewriter strike-out}{ega25_25_122-v0012} linear system of divisors D, one can
canonically associate with it an invertible Module $\underline{L}$ on X, the canonical extension
\var{\inserted[typed]{where U=X-Z.}}{where U=X-Z.}{insertion}{ega25_25_122-v0013}
of $\underline{L}^{x}_{U} = f\ (\underline{O}_P(1))$ , One can give another interpretation of it,
independent of the hypothesis $D\quad\hand{\unsure{\phi|\emptyset}}$ , as follows; by constructio

\var{\xout{the inverse image of $\underline{L}$x}}{}{typewriter strike-out}{ega25_25_122-v0014} or almost $\underline{L}^{x}_{U}\otimes_S\underline{O}_P(1)$ is \var{\xout{none}}{}{typewriter strike-out}{ega25_25_122-v0015} \var{\xout{other than}}{}{typewriter strike-out}{ega25_25_122-v0016} \var{\inserted[typed]{can isomorphic to}}{can isomorphic to}{insertion}{ega25_25_122-v0017} the restriction
$\underline{N}'$ to $U\mathrm{x}_S\overset{\hand{\vee}}{P}$ of the invertible Module $\underline{N}$ on $X\mathrm{x}_S\overset{\hand{\vee}}{P}$ defined by D. By virtue
of \var{\xout{no.1}}{}{typewriter strike-out}{ega25_25_122-v0018} no.1, this determines \var{\xout{uniquely $\underline{N}$, and}}{}{typewriter strike-out}{ega25_25_122-v0019} $\underline{L}_U$ up to unique isom;
and (since U is dense in X) there exists an invertible Module $\underline{L}$ on X,
determined up to unique isom (\unsure{hence} extending $\underline{L}_U$ ) such that there is
an isomorphism $\underline{L}\otimes_{\underline{O}_S}\underline{O}_{\overset{\hand{\unsure{\vee|/}}}{P}}\ (1) \quad \underline{N}$ ; $\underline{L}$ is given for example by
\[ \underline{L} \quad pr_1\ (\underline{N}) \quad ; \]
it is called the invertible Module on X associated with the lin system of divisors
D.
\tnote{25-122-1}{The handwritten slash inside French ``annea/u'' is retained within English ``rin/g''. The final n is absent from ``constructio'' in the source and remains absent in English. ``Hence'' carries the doubt n/ç within French ``donc''. ``Can'', ``isom'', and ``lin'' mean canonically, isomorphism, and linear. The canceled n' and autre que are rendered together as none other than. The missing signs in the displayed expressions remain missing.}

\mspage{122}
